%==============================================================================%
% Abstract
%==============================================================================%

\chapter*{Abstract}
\addcontentsline{toc}{chapter}{Abstract}

\noindent
The transition to post-quantum cryptography demands not only mathematically
sound security arguments but also \emph{machine-verified} correctness proofs for
concrete implementations.  This dissertation presents a comprehensive formal
verification of the Number-Theoretic Transform (NTT) at the core of the
\haetae{} lattice-based digital signature scheme---a Korean national standard
candidate under the KPQC competition.

\medskip
\noindent
We establish functional correctness across three vertically composed
abstraction layers.  At the \emph{implementation layer}, the high-assurance
\jasmin{} programming language produces x86-64 assembly that admits direct
extraction into \easycrypt{} modules.  At the \emph{algorithmic layer}, an
imperative specification in \easycrypt{} faithfully models the in-place
Cooley--Tukey butterfly structure together with Montgomery modular arithmetic
over $\Zq$, $q = 64513$, including tight coefficient-bound invariants across
all eight butterfly stages.  At the \emph{mathematical layer}, a fully abstract
spectral specification expresses the NTT as a Vandermonde-style evaluation at
powers of a primitive $512$th root of unity in $\Zq$.

\medskip
\noindent
The key technical contributions are:
\begin{enumerate}
  \item A machine-checked \textbf{Montgomery reduction correctness} proof,
        showing that $\op{montgomery\_reduce}(a) \equiv a \cdot R^{-1}
        \pmod{q}$ for all $a$ in the $16$-bit to $32$-bit working range, with
        explicit arithmetic bit-width bounds.

  \item A \textbf{butterfly-stage induction} connecting the iterative
        forward NTT loop in \ec{NTT\_Fq.ec} to the closed-form evaluation
        formula in \ec{NTTFullSpec.ec}, parameterised over the 8-level
        Cooley--Tukey recursion for the negacyclic polynomial ring
        $\Z_q[X]/(X^{256}+1)$.

  \item An \textbf{inverse NTT correctness} proof that tracks the final
        normalisation by \(256^{-1}\) and the concrete Montgomery output
        convention, where the exported inverse transform returns coefficients
        represented with one factor of \(R\), as expected for
        \texttt{invntt\_tomont}.

  \item A \textbf{Jasmin-to-EasyCrypt refinement} proof in
        \ec{RefJasminNTT.ec} connecting the extracted Jasmin semantics to the
        imperative specifications, with a 16-bit-to-24-bit coefficient-bound
        analysis for the forward transform and a tighter 16-bit preservation
        proof for the inverse transform.

  \item An \textbf{end-to-end statement} in \ec{NTTEndToEnd.ec} asserting
        that the \jasmin{} exported functions \jasm{poly\_ntt\_jazz} and
        \jasm{poly\_invntt\_jazz} compute exactly the abstract spectral NTT and
        its inverse, through two top-level Hoare statements for the forward
        and inverse exported functions.
\end{enumerate}

\noindent
The verification identifies a critical structural difference between the
\haetae{} NTT and the MLKEM NTT: the former performs a \emph{complete}
$256$-point negacyclic transform with 8-bit bit-reversed twiddle indices,
whereas the latter applies an incomplete $128$-point transform with a
different exponent structure.
This difference invalidates a naive reuse of existing MLKEM algebra lemmas
and motivates new algebraic infrastructure developed in
\ec{NTTFullAlgebra.ec}.

\medskip
\noindent
The scalar NTT verification compiles in \easycrypt{} version~2024 on an
x86-64 Linux platform through the top-level \texttt{NTTEndToEnd.ec} file.
The checked development totals approximately 8\,000 lines of \easycrypt{}
and 245 lines of \jasmin{} support code.

\vspace{1em}
\noindent
\textbf{Keywords:} Number-Theoretic Transform, Formal Verification,
EasyCrypt, Jasmin, Post-Quantum Cryptography, HAETAE, Montgomery Arithmetic,
Lattice-Based Signatures, Functional Correctness, Negacyclic Convolution.
